
\documentclass{article}
\usepackage{colortbl}
\usepackage{makecell}
\usepackage{multirow}
\usepackage{supertabular}

\begin{document}

\newcounter{utterance}

\centering \large Interaction Transcript for game `eqbench', experiment `eq\_bench', episode 6 with reflexion{-}Qwen3.5{-}9B.
\vspace{24pt}

{ \footnotesize  \setcounter{utterance}{1}
\setlength{\tabcolsep}{0pt}
\begin{supertabular}{c@{$\;$}|p{.15\linewidth}@{}p{.15\linewidth}p{.15\linewidth}p{.15\linewidth}p{.15\linewidth}p{.15\linewidth}}
    \# & \multicolumn{2}{c}{Player} && \multicolumn{2}{c}{Game Master} \\
    \hline

    \theutterance \stepcounter{utterance}  
    & & & \multicolumn{4}{p{0.6\linewidth}}{
        \cellcolor[rgb]{0.9,0.9,0.9}{
            \makecell[{{p{\linewidth}}}]{
                \texttt{\tiny{[P1$\langle$GM]}}
                \texttt{You have attempted this same task before and did not succeed. These are the plans you wrote after those attempts.} \\
\texttt{Trial \#1: New plan: I failed because I assigned a high score to "Arrogant" (9) based solely on Gabe's confident dialogue, ignoring the context of the interaction. Ellie's final retort ("I have talent") directly challenged his perceived superiority, which should have triggered a significant drop in his arrogance and a rise in "Disbelief" or "Intimidation" as his ego was bruised. My strategy was too static; I didn't weigh the impact of the final exchange. Next time, I will analyze the *change* in dynamic: Gabe's confidence was met with a direct counter{-}attack, making "Arrogant" the least likely outcome (low score) and "Disbelief" or "Intimidated" the most probable high scores, reflecting the deflation of his ego.} \\
\\ 
\texttt{Trial \#2: New plan: I failed because I assigned a high score to "Arrogant" (3) and "Disbelief" (6) while underestimating the intensity of "Intimidated" (8) and "Respect" (2). The dialogue shows Gabe's confidence crumbling under Ellie's direct challenge ("I don't have time for this," "I have talent"). His final line ("You're going to need it") is not a display of arrogance, but a defensive, fearful warning born of realizing his superiority is an illusion. This specific interaction should trigger a massive spike in "Intimidated" (as his ego is bruised) and "Disbelief" (that his control is lost), while "Arrogant" should be near zero since he is no longer acting superior but rather threatened. "Respect" remains low as he dismisses her, but the primary emotion is fear of her talent. Next time, I will score "Intimidated" as the dominant emotion (9{-}10), "Disbelief" as high (7{-}8), and "Arrogant" as very low (1{-}2), reflecting the shift from confident dismissal to defensive fear.} \\
\\ 
\\ 
\\ 
\texttt{Your task is to predict the likely emotional responses of a character in this dialogue:} \\
\\ 
\texttt{Ellie: So, Gabe, are you ready to lose?} \\
\texttt{Gabe: I don't lose, Ellie. I learn.} \\
\texttt{Ellie: Ah, the eternal student. Must be nice living in a world where failure doesn't exist.} \\
\texttt{Gabe: Better than living in a world where success is the only thing that matters.} \\
\texttt{Ellie: And what's wrong with wanting to succeed? } \\
\texttt{Gabe: Nothing. But when it becomes an obsession, it blinds you.} \\
\texttt{Ellie: And who says I'm obsessed?} \\
\texttt{Gabe: You don't have to say it, Ellie. It's in your eyes, your work, your every breath.} \\
\texttt{Ellie: And what about you, Gabe? What's in your eyes? Fear? Doubt? Or just a whole lot of nothing?} \\
\texttt{Gabe: I guess you'll never know.} \\
\texttt{Ellie: I guess I won't. And maybe I don't want to.} \\
\texttt{Gabe: Maybe you don't. But maybe you do. Maybe you're just too scared to admit it.} \\
\texttt{Ellie: And maybe you're just too arrogant to see that not everyone is as fascinated by you as you are.} \\
\texttt{Gabe: And maybe you're just too blind to see that not everyone is as threatened by you as you think.} \\
\texttt{Ellie: You know what, Gabe? I don't have time for this. I have a masterpiece to create.} \\
\texttt{Gabe: And I have a student to teach. Good luck, Ellie. You're going to need it.} \\
\texttt{Ellie: I don't need luck, Gabe. I have talent.} \\
\texttt{{[}End dialogue{]}} \\
\\ 
\texttt{At the end of this dialogue, Gabe would feel...} \\
\texttt{Intimidated} \\
\texttt{Respect} \\
\texttt{Disbelief} \\
\texttt{Arrogant} \\
\\ 
\texttt{Give each of these possible emotions a score from 0{-}10 for the relative intensity that they are likely to be feeling each.} \\
\\ 
\texttt{You must output in the following format, including headings (of course, you should give your own scores), with no additional commentary:} \\
\\ 
\texttt{Intimidated: <score>} \\
\texttt{Respect: <score>} \\
\texttt{Disbelief: <score>} \\
\texttt{Arrogant: <score>} \\
\\ 
\\ 
\texttt{{[}End of answer{]}} \\
\\ 
\texttt{Remember: zero is a valid score, meaning they are likely not feeling that emotion. You must score at least one emotion > 0.} \\
\\ 
\texttt{Your answer:} \\
            }
        }
    }
     \\ \\

    \theutterance \stepcounter{utterance}  
    & \multicolumn{4}{p{0.6\linewidth}}{
        \cellcolor[rgb]{0.561,0.737,0.561}{
            \makecell[{{p{\linewidth}}}]{
                \texttt{\tiny{[P1$\rangle$GM]}}
                \texttt{Intimidated: 9} \\
\texttt{Respect: 2} \\
\texttt{Disbelief: 7} \\
\texttt{Arrogant: 1} \\
            }
        }
    }
    & & \\ \\

    \theutterance \stepcounter{utterance}  
    & & & \multicolumn{2}{p{0.3\linewidth}}{
        \cellcolor[rgb]{0.9,0.9,0.9}{
            \makecell[{{p{\linewidth}}}]{
                \texttt{\tiny{[GM$|$GM]}}
                \texttt{Intimidated: 9} \\
\texttt{Respect: 2} \\
\texttt{Disbelief: 7} \\
\texttt{Arrogant: 1} \\
            }
        }
    }
    & & \\ \\

    \theutterance \stepcounter{utterance}  
    & & & \multicolumn{2}{p{0.3\linewidth}}{
        \cellcolor[rgb]{0.9,0.9,0.9}{
            \makecell[{{p{\linewidth}}}]{
                \texttt{\tiny{[GM$|$GM]}}
                \texttt{Intimidated: 0} \\
\texttt{Respect: 0} \\
\texttt{Disbelief: 3} \\
\texttt{Arrogant: 4} \\
            }
        }
    }
    & & \\ \\

    \theutterance \stepcounter{utterance}  
    & & & \multicolumn{2}{p{0.3\linewidth}}{
        \cellcolor[rgb]{0.9,0.9,0.9}{
            \makecell[{{p{\linewidth}}}]{
                \texttt{\tiny{[GM$|$GM]}}
                \texttt{game\_result = LOSE} \\
            }
        }
    }
    & & \\ \\

    \theutterance \stepcounter{utterance}  
    & & \multicolumn{4}{p{0.6\linewidth}}{
        \cellcolor[rgb]{0.95,0.95,0.95}{
            \makecell[{{p{\linewidth}}}]{
                \texttt{\tiny{[P1$\leftrightarrow$P]}}
                \texttt{Post{-}episode Reflexion (not sent to GM):} \\
\texttt{New plan: I failed because I assigned a high score to "Arrogant" (1) and underestimated the defensive nature of Gabe's final line, interpreting it as a sign of lingering arrogance rather than a reaction to being threatened. Ellie's statement that Gabe is "too blind to see that not everyone is as threatened by you as you think" directly attacks his self{-}perception, causing his arrogance to deflate into fear. His response, "You're going to need it," is not a confident dismissal but a desperate, fearful warning born of realizing his superiority is an illusion. Therefore, "Arrogant" should be the lowest score (0{-}1), while "Intimidated" should be the dominant emotion (9{-}10) as he feels the weight of her talent and his own vulnerability. "Disbelief" remains high (7{-}8) as he struggles to accept that his control is lost, and "Respect" stays low (1{-}2) as he dismisses her despite his fear. Next time, I will score "Intimidated" as the primary emotion (9), "Disbelief" as secondary (7), "Respect" as minimal (1), and "Arrogant" as near zero (0), reflecting the complete shift from confident superiority to defensive fear.} \\
            }
        }
    }
    & \\ \\

\end{supertabular}
}

\end{document}
